\documentclass[11pt]{article}
\usepackage[T1]{fontenc}
\usepackage{lmodern,microtype,amsmath,amssymb,amsthm,booktabs}
\microtypesetup{expansion=false}
\usepackage[margin=1.05in,headheight=14pt]{geometry}
\usepackage{xcolor}
\definecolor{linknavy}{RGB}{25,55,125}
\usepackage[colorlinks=true,linkcolor=linknavy,citecolor=linknavy,urlcolor=linknavy]{hyperref}
\usepackage{orcidlink,fancyhdr}
\hypersetup{
  pdftitle={Levinson's method in short intervals and simple zeros of the zeta function},
  pdfauthor={Felipe Santibanez-Leal},
  pdfsubject={A localized Levinson-Conrey detector of distinct sign changes and its consequence for simple critical zeros of the Riemann zeta function in short intervals},
  pdfkeywords={Riemann zeta function, simple zeros, critical line, short intervals, Levinson method, mollifier, interval arithmetic}
}
\newtheorem{theorem}{Theorem}[section]
\newtheorem{proposition}[theorem]{Proposition}
\newtheorem{lemma}[theorem]{Lemma}
\newtheorem{corollary}[theorem]{Corollary}
\theoremstyle{remark}
\newtheorem{remark}[theorem]{Remark}
\newcommand{\R}{\mathbb R}
\newcommand{\C}{\mathbb C}
\newcommand{\Z}{\mathbb Z}
\newcommand{\Vt}{\widetilde V}
\newcommand{\wh}{\widehat w(0)}
\newcommand{\doctype}{Preprint}
\newcommand{\docversion}{v0.01}
\newcommand{\docdate}{26 September 2026}
\newcommand{\versiondoi}{10.5281/zenodo.22984155}
\newcommand{\conceptdoi}{10.5281/zenodo.22984154}
\newcommand{\docrepo}{github.com/fsantibanezleal/CAOS\_RESEARCH}
\newcommand{\headerblock}{%
\begin{center}
\fbox{\parbox{0.94\textwidth}{\small\raggedright
\textbf{\doctype} \quad$\cdot$\quad \docversion, \docdate \quad$\cdot$\quad Licence CC BY 4.0\\[2pt]
CAOS Research open-problems programme \quad$\cdot$\quad problem \texttt{riemann-hypothesis}\\[2pt]
Version DOI \href{https://doi.org/\versiondoi}{\texttt{\versiondoi}}\\[2pt]
Concept DOI (always latest) \href{https://doi.org/\conceptdoi}{\texttt{\conceptdoi}}\\[2pt]
Not peer reviewed. Source code, data and computational records: \href{https://\docrepo}{\texttt{\docrepo}}
}}
\end{center}
\vspace{4pt}}
\pagestyle{fancy}
\fancyhf{}
\fancyhead[L]{\small CAOS Research \doctype}
\fancyhead[R]{\small Levinson's method in short intervals\textperiodcentered\ \docversion}
\fancyfoot[C]{\thepage}
\fancypagestyle{plain}{\fancyhf{}\fancyfoot[C]{\thepage}\renewcommand{\headrulewidth}{0pt}}
\allowdisplaybreaks[1]

\title{Levinson's method in short intervals and\\simple zeros of the zeta function}
\author{Felipe Santib\'a\~nez-Leal\,\orcidlink{0000-0002-0150-3246}\\
\small CAOS Research program\\
\small\texttt{github.com/fsantibanezleal/CAOS\_RESEARCH}}
\date{\docdate}

\begin{document}
\maketitle
\thispagestyle{plain}
\headerblock

\begin{abstract}
Fix $1/2<\theta<1$. We localize Levinson's method, with Conrey's general
operator polynomial $Q$, to the intervals $(T,T+T^\theta]$ for mollifiers of
length $T^\nu$ with $\nu<\theta-1/2$, by running Young's short proof of the
mollified second moment with a weight supported on the window. We then show
that on such windows the method counts distinct sign changes of Hardy's
function, that is, distinct critical zeros of odd order, with lower density
$\kappa=1-R^{-1}\log c(P,Q,R,\nu)$ for every $Q$ with $Q(x)+Q(1-x)$ constant.
For operator polynomials of degree $201$ we certify $\kappa>0.7170\,\nu$ for
$\nu$ between $0.0199$ and $0.0999$. Used as the odd-support input of the
Hilbert--parity inequality $(Q-S)(N-O)\ge2(N-S)^2$, together with Wang's
short-interval pair-correlation theorem, this gives a positive proportion of
simple zeros on the critical line in $(T,T+T^\theta]$ for every fixed
$\theta\in[0.534,1)$. The previous exponents for this statement were $0.552$
(Steuding), $0.55019\ldots$ (Wang) and $0.5458838$ (the same parity inequality
with a rank-six Selberg detector). At $\theta=0.5459$ the lower density
exceeds $0.017763$, about one thousand times the previous bound, and
numerically the new bound improves the previous ones for $\theta$ below about
$0.567$.
All constants are certified with exact rational and ball arithmetic and
replayed independently by validated quadrature. The results are asymptotic,
use Wang's recent preprint for the pair term, and do not prove the Riemann
hypothesis.
\end{abstract}

\section{Introduction}

Let $\rho=\beta+i\gamma$ denote the nontrivial zeros of the Riemann zeta
function. For $T,H>0$ let $N(T,H)$ count the zeros with $T<\gamma\le T+H$,
with multiplicity; let $S(T,H)$ count the simple zeros among them with
$\beta=1/2$; and let $O(T,H)$ count the distinct points $t\in(T,T+H]$ at which
Hardy's function $Z(t)=e^{i\vartheta(t)}\zeta(1/2+it)$ changes sign, that is,
the distinct critical zeros of odd multiplicity. Throughout, $H=T^\theta$ with
a fixed $\theta\in(1/2,1)$ and $L=\log T$.

Selberg proved that $O(T,T^\theta)\gg N(T,T^\theta)$ for every
$\theta>1/2$, and Karatsuba extended this to $\theta>27/82$~\cite{Karatsuba}.
These theorems do not give simple zeros. For simple critical zeros, Steuding
localized Levinson's method with a linear operator polynomial, proved a
short-window mollified moment with error $O(T^{1/3+\epsilon}M^{4/3})$, and
obtained a positive proportion for $\theta\ge0.552$~\cite{Steuding,SteudingThesis};
he also observed that the mollifier range available from the classical
off-diagonal treatment gives only $\theta\ge0.693$ for linear $Q$. Wang's short-interval
pair-correlation theorem~\cite{WangShort} gives the lower density
\begin{equation}\label{eq:wangc}
 c(\theta)=2-\frac\theta2-\frac1{\sqrt2}\cot\frac{\theta}{\sqrt2},
\end{equation}
which is positive for $\theta>\theta_0=0.550193964744\ldots$. The finite
Hilbert--parity inequality of~\cite[Proposition 2.5]{CAOSShort} converts any
lower density $k$ for $O/N$ into the simple-critical density
\begin{equation}\label{eq:h}
 h(\theta;k)=\frac{3+k-\sqrt{(1-k)(9-k-8c(\theta))}}4 .
\end{equation}
With the rank-six Selberg detector of Pearce-Crump~\cite{PearceCrump},
$k_6(\theta)=(\theta-1/2)/(4eC_6)$ with $C_6=0.65663\ldots$, the resulting
term is positive exactly for $\theta>\theta_6$, where
$0.5458837<\theta_6<0.5458838$~\cite[Theorem 1.2]{CAOSShort}. Since the
positivity condition is $k>1-2/(2-c(\theta))$, the onset is governed by the
size of the odd-support density near $\theta=1/2$, where $k_6$ has slope only
$1/(4eC_6)\approx0.140$.

This paper replaces the Selberg detector by Levinson's. The mollified second
moment behind Levinson's method is available in short windows (Theorem
\ref{thm:moment}); the localized method counts distinct sign changes of $Z$
for Conrey's operator polynomials of every degree (Theorem \ref{thm:signs});
and polynomials of high degree push the density to about $0.717(\theta-1/2)$,
five times the Selberg slope. The labels [D] and [MV] mark deductions and
finite machine-verified statements.

\begin{theorem}[Simple critical zeros from $T^{0.534}$, D+MV]
\label{thm:main}
For every fixed $\theta\in[0.534,1)$,
\begin{equation}\label{eq:main}
 \liminf_{T\to\infty}\frac{S(T,T^\theta)}{N(T,T^\theta)}
 \ge\max\left\{0,\;c(\theta),\;\frac{c(\theta)+2\kappa_\theta}3,\;
 h(\theta;\kappa_\theta)\right\}>0,
\end{equation}
where $\kappa_\theta$ is any certified value of Table~\ref{tab:kappa} with
$\nu<\theta-1/2$. In particular the right side is at least
$h(\theta;\kappa_{0.0339})$ with $\kappa_{0.0339}>0.0243176419$, which is
positive and increasing on $[0.534,1)$, and with $\nu=\theta-1/2-10^{-4}$
\[
\begin{array}{c|cccccc}
\theta&0.534&0.535&0.54&0.5459&0.55\\\hline
h(\theta;\kappa_\theta)>&1.4806\cdot10^{-5}&0.0015246&0.0090231&0.0177638&0.0237708
\end{array}
\]
\end{theorem}

At every point of the table $c(\theta)<0$ and $h(\theta;\kappa_\theta)$ is the
largest term. Numerically, $h$ with $\kappa\approx0.7173(\theta-1/2)$ stays
above $c(\theta)$ up to $\theta\approx0.567$; beyond that point Wang's term
$c(\theta)$ is the largest in \eqref{eq:main}. The two analytic inputs are the
following.

\begin{theorem}[Mollified moment on a short window, D]\label{thm:moment}
Let $0<\nu<\theta-1/2$, $R>0$, let $P$ be a real polynomial with $P(0)=0$,
$P(1)=1$, and let $Q$ be a real polynomial with $Q(0)=1$. Put
$\sigma_0=1/2-R/L$, $y=T^\nu$, $\Delta=H/L$,
\[
 V(s)=Q\Bigl(-\frac1L\frac{d}{ds}\Bigr)\zeta(s),\qquad
 \psi(s)=\sum_{h\le y}\frac{\mu(h)}{h^{s+1/2-\sigma_0}}
 P\Bigl(\frac{\log(y/h)}{\log y}\Bigr).
\]
Let $w$ be smooth with $0\le w\le1$, $w=1$ on $[T,T+H]$, support in
$[T-\Delta,T+H+\Delta]$ and $w^{(j)}\ll_j\Delta^{-j}$. Then
\[
 \int_\R w(t)\,|V\psi(\sigma_0+it)|^2\,dt=c(P,Q,R,\nu)\,\wh+O(H/L),
\]
\[
 c(P,Q,R,\nu)=1+\frac1\nu\int_0^1\!\!\int_0^1
 \bigl(w_R(v)P'(u)+\nu w_R'(v)P(u)\bigr)^2du\,dv,\qquad w_R(v)=e^{Rv}Q(v).
\]
\end{theorem}

\begin{theorem}[Distinct sign changes, D]\label{thm:signs}
In the setting of Theorem~\ref{thm:moment}, suppose that $Q(x)+Q(1-x)$ is a
nonzero constant. Then
\[
 \liminf_{T\to\infty}\frac{O(T,T^\theta)}{N(T,T^\theta)}\ge
 \kappa(P,Q,R,\nu):=1-\frac1R\log c(P,Q,R,\nu).
\]
\end{theorem}

Theorem~\ref{thm:moment} is Conrey's constant~\cite{Conrey} in the form given
by Young~\cite{Young}, on a window of length $T^\theta$. The only place where
the window length restricts the mollifier is the off-diagonal term, which
forces $\nu<\theta-1/2$, the analogue of Levinson's range $\nu<1/2$ for
$\theta=1$. Theorem~\ref{thm:signs} is a sign-change sharpening of the
classical counting argument: zeros of the Levinson function on the critical
line carry full weight, and the argument of the detector is followed through
its level crossings. The classical sources conclude simple zeros only for
$\deg Q=1$~\cite{Conrey,BCY}; for higher degree the method does not separate
simple zeros from zeros of odd order at least three, and that is exactly the
information the parity inequality requires.

We could not locate a statement of Theorem~\ref{thm:moment} for general $Q$
on windows shorter than $T$, or of Theorem~\ref{thm:signs} for $\deg Q\ge3$;
Conrey's Gaussian-window proposition~\cite{Conrey} applies to windows
$T^{1-\delta}$ with $\delta<1/7-\nu/4$. The search was bounded and is not a
priority determination.

\section{The mollified moment on a short window}\label{sec:moment}

We follow Young~\cite{Young}. His Theorem 2 treats a weight supported in
$[T/4,2T]$ with $w^{(j)}\ll(T/L)^{-j}$ and a mollifier of length $M=T^{\vartheta}$,
$\vartheta<1/2$. Here the weight has support of length at most $2H$, the
derivative scale is $\Delta=H/L$, and $M$ is replaced by $y=T^\nu$. We list
every step in which the length enters; all others are used verbatim. As in
Young's Section 3, it suffices to prove
\begin{equation}\label{eq:Iab}
 I(\alpha,\beta):=\int w(t)\zeta(\tfrac12+\alpha+it)\zeta(\tfrac12+\beta-it)
 |\psi(\sigma_0+it)|^2dt=c(\alpha,\beta)\,\wh+O(H/L)
\end{equation}
uniformly for $\alpha,\beta\ll1/L$, with $c(\alpha,\beta)$ Young's (3.3)
with $\vartheta$ replaced by $\nu$ and $M$ by $y$. We first work on the annuli
$\alpha,\beta\asymp1/L$, $|\alpha+\beta|\gg1/L$ of Young's Lemma 6. There the
factor $G(s)=e^{s^2}p(s)$, $p(s)=((\alpha+\beta)^2-4s^2)/(\alpha+\beta)^2$, of
Young's approximate functional equation (Lemma 4) satisfies
$p(s)\ll L^2(1+|s|^2)$. Hence the weight $V_{\alpha,\beta}(x,t)$ and the
factor $X_{\alpha,\beta,t}$ of that lemma obey, for $t\ge T/2$,
\begin{equation}\label{eq:Vbound}
 t^j\partial_t^jV_{\alpha,\beta}(x,t)\ll_{A,j}L^2(1+x/t)^{-A},\qquad
 X_{\alpha,\beta,t}=(t/2\pi)^{-\alpha-\beta}(1+O(t^{-1})).
\end{equation}
The first bound is Young's (4.3), whose decay factor is printed as
$(1+|t/x|)^{-A}$; the Mellin representation (4.1) gives $(1+x/t)^{-A}$,
which is the form used in his proof of Lemma 5. The Stirling expansion (4.2)
of $g_{\alpha,\beta}(s,t)$ is uniform only for $|s|=o(t^{1/2})$; on every
vertical line used below, the part $|\Im s|\ge T^{1/3}$ is negligible because
$|G(s)|\ll L^2e^{-(\Im s)^2/2}$ there, while
$|g_{\alpha,\beta}(s,t)|\ll(t+|s|)^{O(1)}$.

\begin{lemma}[Twisted integral]\label{lem:twisted}
Let $h,k\le y$ and let $\alpha,\beta$ lie on the annuli. Then
\begin{multline*}
 \int w(t)\Bigl(\frac hk\Bigr)^{-it}\zeta(\tfrac12+\alpha+it)\zeta(\tfrac12+\beta-it)dt
 =\sum_{hm=kn}\frac{\int V_{\alpha,\beta}(mn,t)w(t)dt}{m^{1/2+\alpha}n^{1/2+\beta}}\\
 +\sum_{hm=kn}\frac{\int V_{-\beta,-\alpha}(mn,t)X_{\alpha,\beta,t}w(t)dt}{m^{1/2-\beta}n^{1/2-\alpha}}
 +O_A(T^{-A}).
\end{multline*}
\end{lemma}

\begin{proof}
Insert the approximate functional equation. Its error term contributes
$O(HT^{-A})$ and the terms with $hm=kn$ are the main terms. For $hm\ne kn$,
the bounds \eqref{eq:Vbound} and $w^{(j)}\ll\Delta^{-j}$ give
$\partial_t^j[w(t)V_{\alpha,\beta}(x,t)]\ll L^2(1+x/T)^{-A}\Delta^{-j}$ on the
support, so $j$ integrations by parts give
\[
 \int w(t)\Bigl(\frac{hm}{kn}\Bigr)^{-it}V_{\alpha,\beta}(mn,t)dt
 \ll HL^2\frac{(1+mn/T)^{-A}}{(\Delta|\log(hm/kn)|)^j}
 \le HL^2(1+mn/T)^{-A}\Bigl(\frac{2\sqrt{hkmn}}{\Delta}\Bigr)^j,
\]
since $|\log(hm/kn)|\ge1/(2\sqrt{hkmn})$ for $hm\ne kn$. Put
$\eta=\theta-1/2-\nu>0$ and $\epsilon=\eta/2$. If $mn\le T^{1+\epsilon}$ then
$2\sqrt{hkmn}/\Delta\le2yT^{(1+\epsilon)/2}L/T^\theta=2LT^{-3\eta/4}$, and there
are at most $T^{2}$ such pairs $(m,n)$, so these terms contribute
$O(T^{-A})$ once $j>4(A+5)/(3\eta)$. If $mn>T^{1+\epsilon}$ we do not integrate
by parts: the trivial bound $HL^2(mn/T)^{-A'}$ summed against
$(mn)^{-1/2+O(1/L)}$ gives $O(HL^3T^{(1+\epsilon)/2-\epsilon A'+o(1)})$, which is
$O(T^{-A})$ for $A'\ge(A+3)/\epsilon$. The reflected part is identical because
$t^j\partial_t^jX_{\alpha,\beta,t}\ll|X_{\alpha,\beta,t}|\ll1$.
\end{proof}

Inserting the mollifier, $I=I_1+I_2+O(T^{-A})$, where $I_1$ is Young's (5.1)
with $M$ replaced by $y$ and $I_2$ is $I_1(-\beta,-\alpha)$ with
$X_{\alpha,\beta,t}$ inserted in the $t$-integral. On the support of $w$ we have
$t=T(1+O(H/T))$, so by \eqref{eq:Vbound}
$X_{\alpha,\beta,t}=(T/2\pi)^{-\alpha-\beta}(1+\rho(t))$ with $\rho\ll H/T$.
Write $h=gh'$, $k=gk'$ with $(h',k')=1$; then $hm=kn$ means $m=k'r$, $n=h'r$,
and
\[
 \sum_{h,k\le y}\frac1{\sqrt{hk}}\sum_{hm=kn}
 \frac{(1+mn/T)^{-A}}{(mn)^{1/2-O(1/L)}}
 \ll L\sum_{g\le y}\frac1g\Bigl(\sum_{h'\le y}\frac1{h'}\Bigr)^2\ll L^4.
\]
Bounding with absolute values, the contribution of $\rho$ to $I_2$ is
$\ll H(H/T)L^6=o(H/L)$ because $\theta<1$. Since $(2\pi)^{\alpha+\beta}=1+O(1/L)$
and $|I_1(-\beta,-\alpha)|\ll H$ by the next lemma (the region of the annuli
is invariant under $(\alpha,\beta)\mapsto(-\beta,-\alpha)$),
\begin{equation}\label{eq:I1I2}
 I(\alpha,\beta)=I_1(\alpha,\beta)+T^{-\alpha-\beta}I_1(-\beta,-\alpha)+O(H/L).
\end{equation}

\begin{lemma}[Diagonal]\label{lem:diagonal}
Uniformly on the annuli, $I_1(\alpha,\beta)=c_1(\alpha,\beta)\wh+O(H/L)$ with
\[
 c_1(\alpha,\beta)=\frac1{(\alpha+\beta)\log y}\int_0^1
 \bigl(P'(u)+\alpha\log y\,P(u)\bigr)\bigl(P'(u)+\beta\log y\,P(u)\bigr)du .
\]
\end{lemma}

\begin{proof}
Follow Young's Section 6. After the Mellin representation (6.1) and the
arithmetic identity (6.2), move the $u,v$-contours to $\Re=\delta$ and the
$s$-contour to $\Re s=-\delta+\epsilon_1$ with a fixed small $\epsilon_1>0$. The only pole crossed is $s=0$,
because $G$ vanishes at the pole $s=-(\alpha+\beta)/2$ of
$\zeta(1+\alpha+\beta+2s)$. The only factor involving $w$ is
$\int w(t)g_{\alpha,\beta}(s,t)dt$, and by Young's (4.2),
$|g_{\alpha,\beta}(s,t)|\ll(t/2\pi)^{-\delta+\epsilon_1}(1+|s|^2)$ on the new
contour for $|\Im s|\le T^{1/3}$. This factor is therefore
$\ll HT^{-\delta+\epsilon_1}(1+|s|^2)$, while $|y^{u+v}|=y^{2\delta}$; the
remaining factors are as in Young. The new contour contributes
$\ll HT^{-\delta(1-2\nu)+\epsilon_1}L^{O(1)}=o(H/L)$ for
$\epsilon_1<\delta(1-2\nu)/2$, since $\nu<1/2$. The residue at $s=0$ has
$g_{\alpha,\beta}(0,t)=G(0)=1$ and equals $\wh$ times the expression in
Young's (6.3). Young's Lemma 7 and Section 7 involve neither $w$ nor $T$ and
give the main term with relative error $O(1/L)$.
\end{proof}

\begin{proof}[Proof of Theorem~\ref{thm:moment}]
As in Young's deduction of his Lemma 3 from Lemma 6,
\[
 c_1(\alpha,\beta)+c_1(-\beta,-\alpha)=\int_0^12PP'=1,\qquad
 \frac{1-T^{-\alpha-\beta}}{(\alpha+\beta)\log y}=\frac1\nu\int_0^1T^{-v(\alpha+\beta)}dv .
\]
With \eqref{eq:I1I2} and Lemma~\ref{lem:diagonal} this gives \eqref{eq:Iab} on
the annuli, and the maximum modulus principle extends it to $\alpha,\beta\ll1/L$
because both sides are holomorphic. By Young's (3.4), the left side of
Theorem~\ref{thm:moment} equals
\[
 Q\Bigl(-\frac1L\frac{\partial}{\partial\alpha}\Bigr)
 Q\Bigl(-\frac1L\frac{\partial}{\partial\beta}\Bigr)I(\alpha,\beta)
 \Big|_{\alpha=\beta=-R/L};
\]
the derivatives are Cauchy integrals over circles of radius $\asymp1/L$ on
which the error is uniform, and the operator applied to $c(\alpha,\beta)$ gives
Young's (1.3) with $\vartheta$ replaced by $\nu$. Since
\[
 \frac{d}{dx}\Bigl[e^{R\nu x}P(x+u)Q(v+\nu x)\Bigr]_{x=0}
 =e^{-Rv}\bigl(w_R(v)P'(u)+\nu w_R'(v)P(u)\bigr),
\]
this is the stated constant.
\end{proof}

\begin{remark}\label{rem:uses}
The hypothesis $\nu<\theta-1/2$ is used only in Lemma~\ref{lem:twisted}, and
the bound $\nu<1/2$ only in Lemma~\ref{lem:diagonal}. The hypothesis
$\theta<1$, through $H=o(T)$, places the support of $w$ in $[T/2,2T]$, so that
$t^{-1}\le\Delta^{-1}$, and bounds $\rho$. The normalization $Q(0)=1$ enters
only through the constant term: the operator maps the term $1$ of
$c(\alpha,\beta)$ to $Q(0)^2$, so for every real polynomial $Q$ the moment is
$\bigl(Q(0)^2+\nu^{-1}\int_0^1\!\int_0^1(w_RP'+\nu w_R'P)^2\bigr)\wh+O(H/L)$,
and in particular $O(H)$.
\end{remark}

\section{Distinct sign changes}\label{sec:signs}

\subsection{An exact identity on the critical line}

Let $\chi(s)=\pi^{s-1/2}\Gamma((1-s)/2)/\Gamma(s/2)$, so that
$\zeta(s)=\chi(s)\zeta(1-s)$, and put $\lambda(s)=-\chi'(s)/\chi(s)$. Write
$D=d/ds$ and $Q(x)=\sum_kq_kx^k$.

\begin{lemma}\label{lem:identity}
For every $k\ge0$, $\chi(s)\zeta^{(k)}(1-s)=(-1)^k(D+\lambda)^k\zeta(s)$, where
$(D+\lambda)^k$ is the $k$-fold composition of $f\mapsto f'+\lambda f$.
Consequently $\chi(s)V(1-s)=\sum_kq_kL^{-k}(D+\lambda)^k\zeta(s)$. If
$Q(x)+Q(1-x)=\beta$ identically, then
\begin{equation}\label{eq:Ebeta}
 E_\beta:=\beta\zeta-V-\chi(s)V(1-s)=\sum_kq_kL^{-k}
 \bigl[(D+L)^k-(D+\lambda)^k\bigr]\zeta(s),
\end{equation}
and with $\Vt=V+E_\beta/2$ and $\omega(t)=e^{i\vartheta(t)}$,
\begin{equation}\label{eq:real}
 \beta Z(t)=2\,\Re\bigl(\omega(t)\Vt(\tfrac12+it)\bigr)\qquad(t\in\R).
\end{equation}
\end{lemma}

\begin{proof}
Put $f(s)=\zeta(1-s)=\chi(1-s)\zeta(s)$ and note $-\chi'(1-s)/\chi(1-s)=
\lambda(s)$ because $\chi(s)\chi(1-s)=1$. Then $f'=\chi(1-s)(D+\lambda)\zeta$ and,
inductively, $f^{(k)}=\chi(1-s)(D+\lambda)^k\zeta$; since
$\zeta^{(k)}(1-s)=(-1)^kf^{(k)}(s)$, the first claim follows, and the second by
linearity. The polynomial identity $Q(-x)+Q(1+x)=\beta$ with $x=D/L$ gives
$\beta\zeta=V+\sum_kq_kL^{-k}(D+L)^k\zeta$, which is \eqref{eq:Ebeta}. On
$s=1/2+it$ we have $1-s=\bar s$, $V(\bar s)=\overline{V(s)}$ and
$\chi(s)=\omega^{-2}$, so $\omega\chi(s)V(1-s)=\overline{\omega V(s)}$ and
$\omega\zeta=Z$. Hence $\beta Z=2\Re(\omega V)+\omega E_\beta$. The left side
and $2\Re(\omega V)$ are real, so $\omega E_\beta$ is real, and
\eqref{eq:real} follows.
\end{proof}

The identity \eqref{eq:real} holds for every real $Q$; the symmetry of $Q$
makes $E_\beta$ small. By Stirling's formula, $\lambda(s)=\log(|t|/2\pi)+
O(1/|t|)$ and $\lambda^{(j)}(s)\ll1/|t|$ for $j\ge1$, uniformly in fixed
vertical strips, so $\lambda-L=O(1)$ for $t\asymp T$. In normal-ordered form
$(D+\lambda)^k=\sum_{m\le k}c_{k,m}D^m$, where $c_{k,m}$ is
$\binom km\lambda^{k-m}$ plus a polynomial in $\lambda,\lambda',\dots$ of
weight $k-m$ whose terms contain a derivative of $\lambda$. Hence the
coefficient of $D^m$ in $(D+L)^k-(D+\lambda)^k$ is $O(L^{k-m-1})$, and
\begin{equation}\label{eq:Esize}
 E_\beta=\sum_{j<\deg Q}\varepsilon_j(s)L^{-j}\zeta^{(j)}(s),\qquad
 \varepsilon_j(s)\ll1/L,
\end{equation}
uniformly for $\sigma_0\le\sigma\le\sigma_1$ and $T/2\le t\le2T$. For $j\ge1$,
$L^{-j}\zeta^{(j)}=(-1)^j(V_{1+x^j}-\zeta)$, where $V_{1+x^j}$ is the $V$ of
the polynomial $1+x^j$, so Theorem~\ref{thm:moment} for the polynomials $1$
and $1+x^j$ and Cauchy--Schwarz give
\begin{equation}\label{eq:Emean}
 \int w\,|\psi E_\beta(\sigma_0+it)|^2dt\ll H/L^2 .
\end{equation}
Published admissible polynomials often have $\beta\ne1$; for instance
Conrey's polynomial of degree seven has $\beta=0.984$. For them the
unnormalized difference $\zeta-V-\chi V(1-s)$ is not small, which is why
$E_\beta$ is normalized by $\beta$.

\subsection{The counting lemma}

\begin{lemma}\label{lem:count}
Let $\sigma_1=\sigma_1(P,Q)$ be fixed and large, let $L=\log T_0$ be fixed with
$T_0$ large, and let $|T-T_0|\le1$, $H\asymp T_0^\theta$, where $T$ and $T+H$
are not ordinates of zeros of $Z$ or of $\Vt$ in $[1/2,\sigma_1]$. Let
$N_{\Vt}$ count, with multiplicity, the zeros of $\Vt$ in
$[1/2,\sigma_1]\times(T,T+H]$, including those on $\sigma=1/2$. Then
\[
 O(T,H)\ge N(T,H)-2N_{\Vt}-O(L).
\]
\end{lemma}

\begin{proof}
Let $t_1<\dots<t_J$ be the zeros of $\Vt(1/2+it)$ in $(T,T+H)$, of orders
$k_1,\dots,k_J$, and $K=\sum k_j$. Write
$\Vt(s)=\prod_j(s-1/2-it_j)^{k_j}F(s)$ with $F$ holomorphic and zero-free on the
segment. On the line $\omega\Vt=i^Kp(t)\omega F$ with
$p(t)=\prod_j(t-t_j)^{k_j}$ real, so by \eqref{eq:real}
$\beta Z=2p\,\Re W_0$, where $W_0=i^K\omega F$ is continuous and zero-free on
$[T,T+H]$.

Let $\phi$ be a continuous argument of $W_0$; assume $\phi(T)<\phi(T+H)$, the
other case being symmetric, and let $c_1<\dots<c_M$ be the points of
$\pi/2+\pi\Z$ in $(\phi(T),\phi(T+H))$, so that $M\ge|\Delta\phi|/\pi-1$. Put
$\tau_0=T$, $\tau_M=T+H$, and for $0<i<M$ let $\tau_i$ be the first point with
$\phi(\tau_i)=c_i+\pi/2$. Then $\tau_0<\dots<\tau_M$ and $\cos\phi(\tau_i)$ alternates in
sign, since $\phi(\tau_0)\in(c_1-\pi,c_1)$, $\phi(\tau_M)\in(c_M,c_M+\pi)$ and
$\phi(\tau_i)=c_i+\pi/2$. After moving the interior $\tau_i$ slightly, none of them
is a $t_j$. As $\operatorname{sgn}Z(\tau_i)=\operatorname{sgn}(\beta)
\operatorname{sgn}p(\tau_i)\operatorname{sgn}\cos\phi(\tau_i)$ and $p$ changes sign
only at the $t_j$, the function $Z$ has opposite signs at the ends of at least
$M-J$ of the disjoint intervals $(\tau_{i-1},\tau_i)$, and each such interval
contains a sign change. Hence $O(T,H)\ge|\Delta\phi|/\pi-1-J$.

Apply the argument principle to $F$ on $[1/2,\sigma_1]\times[T,T+H]$. It is
zero-free on the left side, and its zeros inside are the $N_>$ zeros of $\Vt$
with $\sigma>1/2$. On the bottom, right and top sides
$\arg F=\arg\Vt-\sum_jk_j\arg(s-1/2-it_j)$, and each factor turns by exactly
$+\pi$ along these sides. Therefore
\[
 \Delta\phi=\Delta\vartheta+\Delta_{\rm brt}\arg\Vt-\pi K-2\pi N_> .
\]
The Riemann--von Mangoldt formula gives $\Delta\vartheta/\pi=N(T,H)+O(L)$. On
$\sigma=\sigma_1$ we have $|V-1|\le1/4$ and $|E_\beta|\ll1/L$, so $\Re\Vt>0$ and
that side contributes $O(1)$. On the horizontal sides the Backlund--Jensen
argument gives $O(L)$, because $\Vt=(\beta\zeta+V-\chi V(1-s))/2=
\sum_{j\le\deg Q}a_j(s)\zeta^{(j)}(s)$ with each $a_j$ a polynomial in
$L^{-1}$, $\lambda$ and its derivatives, so $\Vt$ is holomorphic and
polynomially bounded near the window and $\Re\Vt\ge1/2$ at $\sigma_1+iT$.
Since $J\le K$ and $N_{\Vt}=N_>+K$,
\[
 O(T,H)\ge N(T,H)-2(N_>+K)-O(L)=N(T,H)-2N_{\Vt}-O(L).\qedhere
\]
\end{proof}

For a general window we keep $L=\log T$ and move both endpoints inward by less
than one unit to avoid the finitely many exceptional ordinates, including those
of zeros of $\psi\Vt$; each unit interval contains $O(L)$ zeros of $\zeta$ and,
by Jensen's formula, $O(L)$ zeros of $\Vt$ and of $\psi\Vt$ in
$[\sigma_0,\sigma_1]$, so every count changes by $O(L)$. The quantity $J$ in
the proof may be replaced by the number of $t_j$ of odd order.

A double zero of $Z$ is not a sign change and is never counted; it is paid
for in $N_{\Vt}$, as an on-line zero of $\Vt$ when $\deg Q=1$ and, in the
examples of Section~\ref{sec:certify}, as a zero of $\Vt$ to the right of the
line when $\deg Q=3$. If the on-line zeros of $\Vt$ are omitted from
$N_{\Vt}$, the inequality fails (Section~\ref{sec:certify}).

\subsection{Proof of Theorem~\ref{thm:signs}}

Let $F_1=\psi\Vt$. Littlewood's lemma on $[\sigma_0,\sigma_1]\times[T,T+H]$
gives
\[
 2\pi\int_{\sigma_0}^{\sigma_1}n(\sigma)d\sigma=
 \int_T^{T+H}\log|F_1(\sigma_0+it)|dt-\int_T^{T+H}\log|F_1(\sigma_1+it)|dt+O(L),
\]
where $n(\sigma)$ counts the zeros of $F_1$ in the rectangle with real part
greater than $\sigma$, and $O(L)$ bounds the horizontal argument integrals:
$\psi$ and $\Vt$ are polynomially bounded and have real part at least $1/2$ at
$\sigma_1+iT$ and $\sigma_1+i(T+H)$, so the Backlund--Jensen argument bounds
the variation of $\arg\psi$ and of $\arg\Vt$, hence of $\arg F_1$. Each zero counted in $N_{\Vt}$ has real
part at least $1/2=\sigma_0+R/L$, and every other zero contributes
nonnegatively, so the left side is at least $2\pi(R/L)N_{\Vt}$.

On $\sigma=\sigma_1$, $\psi V=1+\sum_{n\ge2}b_nn^{-s}$ with
$|b_n|\ll d(n)(\log n)^{\deg Q}$, so its logarithm is an absolutely
convergent Dirichlet series for $\sigma_1$ large and
$\int\log|\psi V(\sigma_1+it)|dt=O(1)$; moreover $|E_\beta/(2V)|\ll1/L$ there,
which contributes $O(H/L)$. On $\sigma=\sigma_0$, concavity of the logarithm,
$w=1$ on $[T,T+H]$ with $w\ge0$, the inequality
$|a+b|^2\le(1+L^{-1/2})|a|^2+(1+L^{1/2})|b|^2$, Theorem~\ref{thm:moment} and
\eqref{eq:Emean} give
\[
 \int_T^{T+H}\log|F_1(\sigma_0+it)|dt\le\frac H2\log\Bigl(\frac1H\int
 w|\psi\Vt|^2\Bigr)\le\frac H2\log c(P,Q,R,\nu)+O(HL^{-1/2}).
\]
Hence $N_{\Vt}\le(HL/4\pi R)\log c+o(HL)$. Since $N(T,H)=(HL/2\pi)(1+o(1))$,
Lemma~\ref{lem:count} gives $O/N\ge1-R^{-1}\log c-o(1)$. \qed

\begin{remark}
The number $\sigma_1$ and all implied constants depend on $P$, $Q$, $R$ and
$\nu$. For the polynomials of Section~\ref{sec:certify} the monomial
coefficients of $Q$ have size about $10^{120}$ and $R$ is about $28$, so the
constants are very large, and for $\eta=\theta-1/2-\nu=10^{-4}$ the factor
$2LT^{-3\eta/4}$ of Lemma~\ref{lem:twisted} is below one only when $\log T$
exceeds about $1.7\cdot10^5$. The results are purely asymptotic.
\end{remark}

\section{Transfer to simple zeros}\label{sec:transfer}

\begin{theorem}[D]\label{thm:transfer}
If $\liminf O(T,T^\theta)/N(T,T^\theta)\ge k\ge0$, then
$\liminf S(T,T^\theta)/N(T,T^\theta)\ge h(\theta;k)$.
\end{theorem}

\begin{proof}
This is the argument of~\cite[proof of Theorem 1.2]{CAOSShort} with the
odd-support input replaced by $k$. For a real even $\eta\in
C_c^\infty((-\lambda/2,\lambda/2))$, $0<\lambda<\theta$, with $\int\eta^2=1$,
let $Q_T$ be the pair sum of $f=\eta^2$ over the scaled zeros of the window;
it is unrelated to the polynomial $Q$ of the previous sections. With
$q_T,s_T,o_T$ the ratios of $Q_T,S,O$ to $N$, the Hilbert--parity inequality
\cite[Proposition 2.5]{CAOSShort} reads $(q_T-s_T)(1-o_T)\ge2(1-s_T)^2$, and
Wang's theorem~\cite[Theorem 2.2]{WangShort} gives $q_T\to\mathcal C(f)$. If
$O=N$, every zero of the window is simple and critical and $s_T=1$. Otherwise
$1-o_T>0$ and the inequality gives $q_T\ge s_T$; for every $\epsilon>0$ and all
large $T$, $q_T\le\mathcal C(f)+\epsilon$ and $o_T\ge k-\epsilon$, and enlarging
the nonnegative factors gives
$2(1-s_T)^2\le(\mathcal C(f)+\epsilon-s_T)(1-k+\epsilon)$. Hence
$\liminf s_T\ge R(\mathcal C(f),1-k)$ with
$R(A,b)=(4-b-\sqrt{b(b+8(A-1))})/4$, the smaller root. Letting $f$ tend to the
cosine density at fixed $\lambda$, so that $\mathcal C(f)\to2-c(\lambda)$, and
then $\lambda\to\theta$, gives $R(2-c(\theta),1-k)=h(\theta;k)$.
\end{proof}

For fixed $k<1$, $h(\theta;k)$ increases with $c$ and with $k$, since
$\partial_k[(1-k)(9-k-8c)]=-(10-2k-8c)<0$ for $c<1$. As
$c'(\theta)=\frac12\cot^2(\theta/\sqrt2)>0$, if $h(\theta_*;\kappa)>0$ for a
detector with $\nu<\theta_*-1/2$, then the same detector is admissible and
$h(\theta;\kappa)>0$ for every $\theta\in[\theta_*,1)$. The linear branch
$(c+2k)/3$ in \eqref{eq:main} is the transfer of~\cite[Theorem 1.4]{CAOSShort}
with the same input.

\section{Certified detector constants}\label{sec:certify}

We use $Q$ in the Chebyshev family
\[
 Q(y)=1-y+\sum_{j=1}^{100}x_j\bigl(T_{2j+1}(1-2y)-(1-2y)\bigr),
\]
which has degree $201$, $Q(0)=1$, $Q(1)=0$ and $Q(y)+Q(1-y)=1$ for all rational
$x_j$, because $T_{2j+1}$ is odd with $T_{2j+1}(1)=1$,
and $P(x)=S(rx)/S(r)$ with $S$ the Taylor polynomial of $\sinh$ of degree
$13$. The coefficients $x_j$ are rationals with denominator $10^{40}$ obtained
by minimizing $c$ numerically, and $R$, $r$, $\nu$ are rational. With
$B_P=\int_0^1P^2$, $C_P=\int_0^1P'^2$ and $\int_0^1PP'=1/2$, and since
$\int_0^1w_Rw_R'=(e^{2R}Q(1)^2-Q(0)^2)/2=-1/2$,
\[
 c(P,Q,R,\nu)=\frac12+\frac{C_P}\nu\int_0^1e^{2Rv}Q(v)^2dv
 +\nu B_P\int_0^1e^{2Rv}\bigl(RQ(v)+Q'(v)\bigr)^2dv .
\]
Each integral $\int_0^1e^{2Rv}F(v)dv$ of a rational polynomial equals
$e^{2R}a(F)+b(F)$ with rational $a,b$, by the recursion
$\int_0^1e^{tv}v^kdv=e^ta_k+b_k$, $a_k=1/t-(k/t)a_{k-1}$,
$b_k=-(k/t)b_{k-1}$, $a_0=1/t$, $b_0=-1/t$. Hence $c=Ae^{2R}+B$ with exact
rationals $A,B$, and $\kappa=1-R^{-1}\log c$ needs one exponential and one
logarithm in ball arithmetic (Arb, $4000$ bits). The admissibility identities
are verified in exact rational arithmetic.

\begin{table}[ht]
\centering
\caption{Certified detector constants. Lower bounds are rounded down.}
\label{tab:kappa}
\begin{tabular}{llllll}
\toprule
$\nu$ & $R$ & $r$ & $c(P,Q,R,\nu)$ & $\kappa>$ & $\kappa/\nu>$\\
\midrule
$0.0199$ & $48.391$ & $0.7702$ & $5.1998\cdot10^{20}$ & $0.0142729911$ & $0.717235$\\
$0.0299$ & $32.244$ & $0.7708$ & $5.0471\cdot10^{13}$ & $0.0214483010$ & $0.717334$\\
$0.0339$ & $28.439$ & $0.7708$ & $1.1235\cdot10^{12}$ & $0.0243176419$ & $0.717334$\\
$0.0349$ & $27.625$ & $0.7709$ & $4.9777\cdot10^{11}$ & $0.0250349766$ & $0.717334$\\
$0.0399$ & $24.163$ & $0.7709$ & $1.5614\cdot10^{10}$ & $0.0286216496$ & $0.717334$\\
$0.0458$ & $21.05$ & $0.7708$ & $6.9431\cdot10^{8}$ & $0.0328539236$ & $0.717334$\\
$0.0499$ & $19.321$ & $0.7709$ & $1.2321\cdot10^{8}$ & $0.0357949954$ & $0.717334$\\
$0.0999$ & $9.648$ & $0.7707$ & $7.7588\cdot10^{3}$ & $0.0716640051$ & $0.717357$\\
\bottomrule
\end{tabular}
\end{table}

In our numerical optimization the ratio $\kappa/\nu$ did not exceed about
$0.7173$ for this mollifier shape at small $\nu$; the polynomials of Conrey,
Farmer, Kwan, Lin and Turnage-Butterbaugh~\cite{CFKL} give about $0.682\nu$.
Linear $Q$ does not suffice at these lengths: $Q=1-x$, $P=x$ needs
$\nu>0.1928$ for $\kappa>0$. The same code reproduces Young's
$c=2.35006777\ldots$ for $P=x$, $Q=1-x$, $R=1.3$, $\nu=1/2$, and Conrey's
$\kappa>0.4088$ for his polynomial of degree seven with $\nu=4/7$.

An independent program evaluates $Q$ by the Chebyshev recurrence directly
from the $x_j$ and $P$ by its series, separates the double integral by
Fubini into one-dimensional integrals, computes each by validated Arb
quadrature, and recomputes $c(\theta)$ and $h$ with $100$-digit interval
arithmetic in \texttt{mpmath}. Every enclosure overlaps the primary one, and
all thresholds hold again. With $\nu=\theta-1/2-10^{-4}$, the directed values
of $h(\theta;\kappa)$ at $\theta=0.534,0.535,0.54,0.5459,0.55$ exceed
$1.4806994\cdot10^{-5}$, $0.0015246940$, $0.0090231376$, $0.0177638490$ and
$0.0237708528$; at $\theta=0.5459$ the ratio to the rank-six Selberg value
$1.7764518\cdot10^{-5}$ exceeds $999.96$.

Numerical controls support Lemma~\ref{lem:count}. At $T=10^4,10^5,10^6$ and
for $Q$ of degrees $1,3,5,7$, including Conrey's polynomial and polynomials
with $\beta\ne1$, the direct evaluation of $\chi(s)V(1-s)$ agrees with
Lemma~\ref{lem:identity}, $\omega E_\beta$ is real and \eqref{eq:real} holds to
the working precision. On windows at these heights the sign changes of $Z$,
the level crossings of $\arg W_0$ and the zero counts coincide. For the test
functions $f=\zeta\prod_j((s-1/2)^2+\gamma_j^2)^{e_j}$ with planted double
($e_j=1$) and triple ($e_j=2$) critical zeros, the inequality of
Lemma~\ref{lem:count} holds with equality in the double-zero cases, double
zeros are not counted, and the variant that omits the on-line zeros of $\Vt$
fails for $\deg Q=1$.

\section{Scope}

Theorems~\ref{thm:moment} and~\ref{thm:signs} are unconditional and use only
classical tools together with Young's argument. Theorem~\ref{thm:main} also
uses Wang's short-interval pair theorem~\cite{WangShort}, a recent preprint,
through the finite inequality of~\cite{CAOSShort}. All statements are
asymptotic and give no effective starting height. They do not show that
every zero in the window is simple or on the critical line, and they do not
prove the Riemann hypothesis.

The onset is limited by the mollifier range $\nu<\theta-1/2$, which is the
classical range; with linear $Q$ it gives only $\theta\ge0.693$. The gain
comes from operator polynomials of high degree, which lose the simple-zero
information, together with the parity inequality, which restores it. A
mollified moment with general $Q$ in the range $\nu<(3\theta-1)/4$, which
Steuding established for $Q(x)=1-x$~\cite{SteudingThesis}, would, by a
numerical evaluation with slopes as in Table~\ref{tab:kappa}, make the
positivity condition $\kappa>1-2/(2-c(\theta))$ hold for every $\theta>1/2$;
we do not prove such a moment here.

\medskip
\noindent\textbf{Data availability.} The programs, exact parameters, outputs
and independent replays are archived as computation EXP-010 in the
repository~\cite{Record}.

\medskip
\noindent\textbf{Declaration of generative AI use.} Generative AI tools were
used for source research, derivation, review and implementation. The author
reviewed and verified all content and takes full responsibility for the work.

\begin{thebibliography}{99}
\raggedright

\bibitem{BCY}
H. M. Bui, J. B. Conrey and M. P. Young,
\emph{More than 41\% of the zeros of the zeta function are on the critical
line}, Acta Arith. \textbf{150} (2011), 35--64; arXiv:1002.4127.

\bibitem{CAOSShort}
F. Santib\'a\~nez-Leal,
\emph{Simple critical zeros in short intervals: stability, parity,
localization, Hilbert compression, and spectral defect},
Zenodo preprint (2026),
\url{https://doi.org/10.5281/zenodo.22860012}.

\bibitem{CFKL}
J. B. Conrey, D. W. Farmer, C.-H. Kwan, Y. Lin and
C. L. Turnage-Butterbaugh,
\emph{Short mollifiers of the Riemann zeta-function},
arXiv:2508.11108v1 (2025).

\bibitem{Conrey}
J. B. Conrey,
\emph{More than two fifths of the zeros of the Riemann zeta function are on
the critical line}, J. Reine Angew. Math. \textbf{399} (1989), 1--26.

\bibitem{Karatsuba}
A. A. Karatsuba,
\emph{On the zeros of the function $\zeta(s)$ on short intervals of the
critical line}, Math. USSR-Izv. \textbf{24} (1985), 523--537.

\bibitem{PearceCrump}
A. Pearce-Crump,
\emph{Optimising Selberg's method for critical zeros},
arXiv:2609.15329v1 (2026),
\url{https://arxiv.org/abs/2609.15329v1}.

\bibitem{Record}
F. Santib\'a\~nez-Leal,
\emph{EXP-010: localized Levinson parity transfer},
CAOS Research computational record (2026),
\url{https://github.com/fsantibanezleal/CAOS_RESEARCH/tree/main/problems/number-theory/riemann-hypothesis/experiments/EXP-010-levinson-parity-transfer}.

\bibitem{Steuding}
J. Steuding,
\emph{On simple zeros of the Riemann zeta-function in short intervals},
Acta Math. Hungar. \textbf{96} (2002), 259--308.

\bibitem{SteudingThesis}
J. Steuding,
\emph{On simple zeros of the Riemann zeta-function},
doctoral dissertation, Universit\"at Hannover (1999).

\bibitem{WangShort}
B. Wang,
\emph{Simple critical zeros and distinct zeros of the Riemann zeta-function
in short intervals},
arXiv:2609.07918v1 (2026),
\url{https://arxiv.org/abs/2609.07918v1}.

\bibitem{Young}
M. P. Young,
\emph{A short proof of Levinson's theorem},
Arch. Math. \textbf{95} (2010), 539--548; arXiv:1002.4403.

\end{thebibliography}
\end{document}
